\section{TopologyPipeline and SkillPipeline: Agent Topology Optimization and Artifact-Agnostic Skill Detection}
\label{sec:topology_skill}

\subsection{Overview}

The TopologyPipeline and SkillPipeline represent UMAF's meta-analysis layer. The TopologyPipeline implements a meta-cognitive capability: analyzing arbitrary multi-agent tasks, proposing candidate topologies using four canonical patterns, evaluating them across five dimensions, and producing executable specifications. The SkillPipeline implements a fan-out/fan-in graph where a scanner classifies artifacts and four parallel detectors analyze skills across independent dimensions.

\subsection{TopologyPipeline: 4-Node Linear Graph}

The pipeline follows a strict linear workflow with no cycles: \texttt{analyzer $\rightarrow$ designer $\rightarrow$ evaluator $\rightarrow$ writer}.

\subsubsection{TopologyAnalyzerRole: 6-Factor Assessment}

The analyzer assesses input tasks across six orthogonal dimensions: data\_dependencies (inter-agent data flow), parallelism\_opportunities (independent sub-tasks), tool\_requirements (distinct tools per agent), error\_domains (distinct failure modes), latency\_sensitivity (real-time vs. batch), and scale (sub-task count and workload type). Each factor is assigned a level (low/medium/high) with reasoning. The fallback uses word count heuristics with conservative defaults.

\subsubsection{TopologyDesignerRole: 4-Pattern Catalog}

Four canonical multi-agent topologies \cite{gdesigner,argdesigner,mass}:
\begin{enumerate}
    \item \textbf{Sequential}: Fixed-order linear chain. Best for strictly ordered workflows.
    \item \textbf{Fan-Out/Fan-In}: Dispatcher $\rightarrow$ parallel workers $\rightarrow$ collector. Best for embarrassingly parallel sub-tasks.
    \item \textbf{Debate/Consensus}: Multiple independent solvers $\rightarrow$ judge synthesizes. Best for high-stakes decisions.
    \item \textbf{Hierarchical}: Orchestrator $\rightarrow$ domain leads $\rightarrow$ workers. Best for complex multi-layered tasks.
\end{enumerate}
Each topology specifies agent configurations, connections, parallelism strategy, strengths, and weaknesses. The fallback generates three default topologies (sequential, fan-out/fan-in, hierarchical) with pre-configured agents.

\subsubsection{TopologyEvaluatorRole: 5-Dimension Scoring}

Scores each topology 1--10 on: latency (critical path length), reliability (fault tolerance), cost\_efficiency (total tokens/calls), simplicity (implementation ease), and scalability (growth handling). The fallback heuristic encodes deliberate tradeoffs: fan-out/fan-in scores highest (36/50), debate\_consensus lowest (26/50) despite highest reliability.

\subsubsection{TopologyWriterRole: Dual-Format Output}

Produces \texttt{topology\_spec.json} (recommended topology with implementation guide and ASCII flow diagrams) and \texttt{topology\_report.md} (comparison table, recommendation rationale, implementation notes). The writer has only \texttt{write\_file}---the most restricted role in UMAF.

\subsection{SkillPipeline: Fan-Out/Fan-In Graph}

The graph topology directly mirrors the fan-out/fan-in pattern \cite{aflow} that the TopologyPipeline's evaluator scores highest: \texttt{scanner $\rightarrow$ [4 parallel detectors] $\rightarrow$ aggregator $\rightarrow$ writer}.

\subsubsection{SkillScannerRole: 8-Category Classification}

The scanner performs four phases: (1) surface scan via \texttt{find} and \texttt{ls}, (2) artifact classification into 8 types (software\_project, research\_paper, blog\_article, documentation, dataset, design\_document, presentation, configuration) using weighted scoring from multiple evidence sources, (3) deep reading of up to 12 representative files, (4) structure analysis. The \texttt{\_classify\_artifact()} algorithm computes weighted scores from source file ratio, extension matches, build file matches, and keyword matches---operating fully deterministically without LLM involvement.

\subsubsection{Four Artifact-Agnostic Detectors}

Each detector reads the same \texttt{artifact\_analysis.json} but asks fundamentally different questions:

\begin{itemize}
    \item \textbf{DomainExpertiseDetectorRole}: ``What does the creator know deeply about?'' Uses \texttt{\_DOMAIN\_SIGNALS} catalog across 9 pre-defined domains with indicator terms and proficiency heuristics ($\ge$5 matches $\rightarrow$ ``advanced'').

    \item \textbf{TechnicalCraftDetectorRole}: ``How skilled is the creator at the medium itself?'' Adaptive behavior: switches between code craft detection (design patterns, error handling, type systems) and writing craft detection (argumentation, technical writing) based on artifact type.

    \item \textbf{MethodologyDetectorRole}: ``What tools, workflows, and processes are evident?'' \cite{wirl,stride} Covers 35+ tools across categories (languages, testing, CI/CD, containers, documentation, data science). File signal weighting: file matches count double in proficiency calculations.

    \item \textbf{RigorDetectorRole}: ``How thorough, careful, and complete is the work?'' Uses quantitative metrics: test file ratio, documentation ratio, quality config presence. Detects linting/formatting tools and academic rigor indicators (citations, methodology sections, reproducibility).
\end{itemize}

\subsubsection{SkillAggregatorRole and Output}

\raggedright
The aggregator merges four detector reports with name-based deduplication (keep highest proficiency), category normalization (tools mapped via \texttt{\_CATEGORY\_MAP}, skills via \texttt{\_SKILL\_CATEGORY\_MAP}), and summary statistics. The writer produces \texttt{skills.json} (machine-readable) and \texttt{skills\_report.md} (human-readable with proficiency badges and category emojis).

\subsection{Meta-Circular Validation}

Both pipelines were originally generated by UMAF's own CoderPP pipeline---a recursive self-application validating the framework's code generation capabilities. The TopologyPipeline optimizing topologies and the SkillPipeline detecting skills form a closed validation loop: the framework's own analysis tools confirm its architectural choices.
